% ============================================================
%  Objetivo 1 — Dominios en Conserved Domains (NCBI)
% ============================================================

La secuencia que usamos en todo el TP es la \textbf{Proteína X} de \textit{Clostridium tetani},
de 617 aminoácidos. Según los datos del enunciado absorbe luz como una hemoproteína y está
relacionada con quimiotaxis y estrés oxidativo/nitroxativo.

% ============================================================
\section{Ejercicio 1a}

\begin{consigna}
¿Qué clase de proteínas encuentra? ¿Tienen algún grupo prostético? ¿Qué molécula unen y qué
función tienen informadas?
\end{consigna}

Corrimos BLASTp contra la base nr sin filtrar por organismo y con todos los parámetros por
defecto. Los primeros hits fueron:

\begin{table}[H]
\centering
\footnotesize
\begin{tabular}{llP{2cm}P{2cm}}
\toprule
\# & Descripción (accesión) & Organismo & Identidad (\%) \\
\midrule
1 & methyl-accepting chemotaxis protein (AAO34825.1) & \textit{C. tetani} E88 & 98.5 \\
2 & chemotaxis protein (RXI74251.1)                  & \textit{C. tetani}     & 98.2 \\
3 & chemotaxis protein (RXI44462.1)                  & \textit{C. tetani}     & 98.2 \\
4 & chemotaxis protein (AVP55927.1)                  & \textit{C. tetani}     & 98.2 \\
5 & heme NO-binding domain-containing protein (WP\_035125471.1) & \textit{C. tetani} & 98.8 \\
\bottomrule
\end{tabular}
\caption{Primeros hits del BLASTp de la Proteína X contra nr (E-value = 0.0 para todos).}
\end{table}

Todos los hits son \textbf{proteínas de quimiotaxis de tipo MCP} (\textit{Methyl-accepting
Chemotaxis Protein}) con un dominio de unión a hemo. Varias aparecen anotadas directamente
como ``heme NO-binding domain-containing protein'', lo que ya da una pista muy clara de la
función.

Estas proteínas tienen como \textbf{grupo prostético} un \textbf{grupo hemo} (una porfirina
con un átomo de hierro en el centro), que está coordinado axialmente por una histidina proximal
conservada. La molécula que unen es \textbf{óxido nítrico (NO)}, aunque algunas también pueden
unir O$_2$ dependiendo del organismo. Su función es actuar como \textbf{sensores de NO}
acoplados al sistema de quimiotaxis: cuando detectan NO en el entorno, transmiten esa señal al
dominio MCP, que regula el movimiento bacteriano.

% ============================================================
\section{Ejercicio 1b}

\begin{consigna}
¿Las de \textit{C. tetani} están poco o muy conservadas? ¿En qué otras especies aparecen
homólogas? ¿Con qué grado de conservación?
\end{consigna}

Dentro de \textit{C. tetani} la proteína está \textbf{muy conservada}: las distintas cepas
tienen entre 98 y 99\% de identidad, con E-value 0.0 en todos los casos. Lo que llama la
atención es que los 100 hits de la búsqueda pertenecen todos a \textit{Clostridium} o
Clostridiaceae, sin que aparezca ningún organismo fuera del género con blastp estándar.

\begin{table}[H]
\centering
\footnotesize
\begin{tabular}{lP{2.5cm}P{2.5cm}}
\toprule
Organismo & Identidad (\%) & E-value \\
\midrule
\textit{Clostridium tetani} (otras cepas) & 98--99 & 0.0 \\
\textit{Clostridium cochlearium}           & $\sim$83 & 0.0 \\
\textit{Clostridium tetanomorphum}         & $\sim$67 & 0.0 \\
\textit{Clostridium lundense}              & $\sim$66 & 0.0 \\
\textit{Clostridium botulinum}             & $\sim$61 & 0.0 \\
\textit{Clostridium carboxidivorans}       & $\sim$60 & 0.0 \\
\bottomrule
\end{tabular}
\caption{Conservación de la Proteína X en otras especies de \textit{Clostridium}.}
\end{table}

La identidad cae bastante al salir de \textit{C. tetani}: de 98\% baja a $\sim$83\% en
\textit{C. cochlearium} y a $\sim$60--67\% en otras especies del género. Esto sugiere que la
proteína está \textbf{moderadamente conservada} en el género pero que la especificidad de
secuencia es alta dentro de la especie. Homólogos más distantes (como la guanilato ciclasa
soluble humana, que también tiene un dominio HNOB) no aparecen con blastp estándar porque la
identidad cae por debajo del umbral de sensibilidad de BLAST.

% ============================================================
\section{Ejercicio 1c}

\begin{consigna}
¿Cuántos dominios presenta la proteína según el NCBI-CDD? ¿De qué fuente se obtuvo cada uno?
¿Qué funciones tienen descritas?
\end{consigna}

En la solapa \textbf{Graphic Summary} del resultado de BLAST, el NCBI-CDD identifica
\textbf{2 dominios} en la Proteína X:

\begin{table}[H]
\centering
\footnotesize
\begin{tabular}{lllP{5.5cm}}
\toprule
Dominio & Posición & Fuente & Función \\
\midrule
\textbf{HNOB} & $\sim$1--190 aa (N-term.) & CDD/Pfam PF07730 + SMART SM00495 & Unión de grupo hemo; sensor de NO y O$_2$; activa señalización por cambio conformacional \\[4pt]
\textbf{Tar} (MCP) & $\sim$270--617 aa (C-term.) & CDD/Pfam PF00672 + COG0840 & Transducción de señal en quimiotaxis; incluye sitio de dimerización y unión a CheW \\
\bottomrule
\end{tabular}
\caption{Dominios identificados por NCBI-CDD en la Proteína X.}
\end{table}

Además de los dos dominios, el Graphic Summary también marca dos sitios funcionales dentro del
dominio Tar: la \textbf{dimer interface} y el \textbf{putative CheW interface}, que son las
regiones de contacto con otras subunidades del receptor y con la proteína CheW del sistema de
quimiotaxis.

Vale aclarar que NCBI-CDD llama ``Tar'' al dominio C-terminal (por el receptor de aspartato
prototípico de \textit{E. coli}), pero en Pfam el mismo dominio se llama ``MCPsignal'': son
la misma familia con distinto nombre según la base de datos.
